%
%============= CHAPTER BREAK =================%
%
\chapter{Marcus Aurelius (172--180 CE)}%
\vspace*{3mm}%
\justifying%
%
\noindent{}Marcus Aurelius, famed as a philosopher emperor, %
prided himself on his capacity to read and compose Greek. %
His best known work was entitled Τὰ εἰς ἑαυτόν / \textit{Ad se ipsum} / \textit{To Himself}, %
commonly known in English as the \textit{Meditations}. %
Originally written as remembrances, or notes to himself, %
he may not have intended them for publication and distribution. %
There are some clear indications and faint clues scattered among the twelve books of the \textit{Meditations} %
as to when some of them were composed. %
Others lack such references. %
Generally, all of them fall between the years 172 and 180 CE, %
spanning his military campaign along the Danube and approaching the end of his life. \\[1.5mm] %

\noindent{}\textbf{περπερεύομαι} / \textit{perpereuomai} / “to flaunt oneself” %
is fairly rare in the TLG, appearing some 117 times. %
Its first non-Christian occurrence is in book 5 of the \textit{Meditations} of Marcus Aurelius. %
While book 5 has some of the most personal reflections, %
it nevertheless lacks the sort of specific details that might supply a precise date. %
It is safe to place its composition in modest date range, between 173 and 178 CE. %

\vspace*{2mm} 
 \begin{quote}\textls[-15]{οὐκ αἰσθάνῃ πόσα ἤδη παρέχεσθαι δυνάμενος, %
 ἐφ’ ὧν οὐδεμία ἀφυίας καὶ ἀνεπιτηδειότητος πρόφασις, %
 ὅμως ἔτι κάτω μένεις ἑκών; %
 ἢ καὶ γογγύζειν καὶ γλισχρεύεσθαι καὶ κολακεύειν καὶ τὸ σωμάτιον καταιτιᾶσθαι %
 καὶ ἀρεσκεύεσθαι καὶ \textbf{περπερεύεσθαι} καὶ τοσαῦτα ῥιπ\-τάζεσθαι %
 τῇ ψυχῇ διὰ τὸ ἀφυῶς κατεσκευάσθαι ἀναγκάζῃ; (\textit{Medit}. 5.5.1)}\\[1mm] %
 Do you not perceive all that you are capable of providing from your own means, %
 for which there is not any excuse of a lack of natural power or unfitness, %
 nevertheless you still remain below purposely? %
 Or are you compelled to grumble and to be stingy and to flatter and to blame the poor body %
 and to be complaisant and \textbf{to flaunt yourself} and in such ways to be tossed about in soul %
 on account of being uncleverly prepared? \end{quote} % 
\vspace*{2mm}

The first occurrence of the lemma in a Christian work is in the famous love chapter %
of 1 Corinthians, specifically in 1Cor 13:4. %
While this chapter is attested for the Marcionite version, %
our teamʼs reconstruction included 1Cor 13:4a but omitted 1Cor 13:4b, where this word occurs. % 

\vspace*{2mm} 
 \begin{quote}Ἡ ἀγάπη μακροθυμεῖ, χρηστεύεται ἡ ἀγάπη, %
 οὐ ζηλοῖ, οὐ \textbf{περπερεύεται}, οὐ φυσιοῦται, %
 ⁵ οὐκ ἀσχημονεῖ, οὐ ζητεῖ τὰ ἑαυτῆς, %
 οὐ παροξύνεται, οὐ λογίζεται τὸ κακόν, %
 ⁶ οὐ χαίρει ἐπὶ τῇ ἀδικίᾳ, συγχαίρει δὲ τῇ ἀληθείᾳ· (1Cor 13:4--6) \\[1mm] %
 Love suffers long; love shows kindness, %
 envies not, \textbf{flaunts itself} not, puffs itself up not, %
 ⁵ does not act improperly, does not seek the things of oneself, %
 is not provoked, does not count up the evil, %
 ⁶ does not celebrate over wickedness, but co-celebrates in the truth. %
 \end{quote} % 
\vspace*{2mm}

περπερεύομαι / \textit{perpereuomai} has a middle voice and reflexive sense, wherein the self is subject and object. %
In the sentence in the \textit{Meditations} where this lemma appears, %
it is one of six middle-voice verbs within a sentence of 25 words, %
an unusually high ratio in Greek prose. %
The gratuitous use of reflexive verbs befits a modality wherein the author speaks to himself to chasten himself. %
The Marcionite \textit{Apostolos} tends to favor simple, active verbs, %
while middle-voice verbs occur at far higher rates in the canonical Paulines. %
Along with the broader tendency in canonical 1-2 Corinthians and Romans %
for the redactor to use self-implicating language, %
here we may be witnessing one of its sources of inspiration: %
the self-chastening philosophical reflections of the emperor himself.\endnote{%
1Cor 13:3 should probably be omitted in a future edition of the Marcionite reconstruction, %
based on the prevalence of middle voice verbs, %
self-implicating language, and the theme of voluntary martyrdom, % 
which was typical of Ignatius and early Christian martyr acts.} % %

This lemma does not seem to appear in any digitized Greek papyri or Greek inscriptions. %
After 1Cor 3:4, the lemma next appears in Christian writings in the erudite Clement of Alexandria, %
who quotes the Pauline verse twice (\textit{Paed}. 3.1.3.1, \textit{QDS} 38.2), %
but who also uses it once without reference to that text (\textit{Paed}. 3.11.77.1). %
The cognate adjective πέρπερος / \textit{perperos} / “vainglorious|braggart” first %
appears in the TLG in Polybiusʼ \textit{Histories} (32.2.5, 39.1.3), written in the 2nd century BCE, %
an adjective Epictetus seems to have used once (\textit{Diss}. 3.2.15, 108 CE), %
before Herodian included it in his \textit{General Prosody}.\endnote{\textit{GG} 3.1:194.} %
Clement of Alexandria is the first in the TLG to use the cognate noun %
περπερεία / \textit{perpereia} / “vainglory” (\textit{Paed}. 3.1.3). %
The broader patterns of reception favor περπερεύομαι / \textit{perpereuomai} as a neologism %
that first emerged within, or at least around the time of, the \textit{Meditations} of Marcus Aurelius, %
before it entered into Christian literature via the canonical redaction to the Pauline letters. %
